\section*{Exercises}
\addcontentsline{toc}{section}{Exercises}

The solutions to the exercises that do not ask for code are in
\hyperref[ch:losningar]{appendix~A}. The code is published in the course repository after the
lecture.

\exercise{Predict}{Calculation}
\label{ex:9:1}
Answer before you compile.
\begin{enumerate}[a)]
  \item What is the command byte for a write to \code{RegTxDataHi}? For a read of the same?
  \item Which sequence of five bytes does \code{writeReg(RegTxDataLo, 0xAABB0000U)} produce?
  \item What does \code{readReg()} return if the transport answers \code{FF 12 34 56 78} to the
    five \code{transfer()} calls?
  \item What does it return if you accidentally count all five received bytes?
\end{enumerate}

\exercise{The sign extension}{Code}
\label{ex:9:2}
Write \code{readReg()} \textbf{wrongly} once, with the shift before the cast. Let the transport
answer so that the first data byte is \code{0x80}.
\begin{enumerate}[a)]
  \item What does \code{value} become, and why?
  \item For which values does the incorrect code work anyway?
  \item What would the symptom look like if the bug were still there at bring-up? Which \code{id}
    values would have worked?
\end{enumerate}

\exercise{A transport to debug with}{Code}
\label{ex:9:3}
Write a \code{PrintTransport} that implements \code{ByteTransport} and prints every call:
\begin{console}
select
  -> 81   <- 00
  -> 00   <- 00
  -> 00   <- 00
  -> 01   <- 00
  -> 23   <- 00
deselect
\end{console}
Compare the output line by line against the example in \kapref{ch:spi}. It is not a test double
but a debugging tool, and it is worth keeping for the whole project.

\exercise{The layer boundary}{Reasoning}
\label{ex:9:4}
True or false, with a justification:
\begin{enumerate}[a)]
  \item \code{byte_transport.hpp} may include \code{frame.hpp}.
  \item \code{ByteTransport} ought to have a method \code{writeRegister(index, value)}.
  \item \code{ByteTransport} ought to do the whole five-byte transaction in one call.
  \item \code{readReg()} ought to be public so that the tests can call it directly.
\end{enumerate}
Statement c) sounds the most reasonable and is wrong anyway. Why?
